\documentclass[11pt]{article}
\usepackage[margin=1in]{geometry}
\usepackage{amsmath,amssymb,amsthm,mathtools}
\usepackage[hidelinks]{hyperref}
\newtheorem{theorem}{Theorem}
\newtheorem{lemma}[theorem]{Lemma}
\newtheorem{corollary}[theorem]{Corollary}
\newcommand{\C}{\mathbb C}
\newcommand{\Q}{\mathbb Q}
\newcommand{\Qb}{\overline{\mathbb Q}}
\newcommand{\E}{\mathbb E}
\newcommand{\fall}[2]{#1^{\underline{#2}}}
\title{Coupled angular products and Gamma block-radial limits}
\author{Internal proofs and exact certificates; no priority claim}
\date{16 November 2026}
\begin{document}
\maketitle
\noindent\textbf{Progress and boundary.}
The argument for higher-dimensional FC checks the complete block-radial limit of
\href{sources/long-eit-vi-current.tex}{Long--Le Blanc, EIT VI,
Theorem ``Block-radial limit''}, extends it to every positive real Gamma
shape and polynomial markers, allows additional inactive variables with
arbitrary strictly lower-degree couplings, and proves a genuinely coupled angular
nonvanishing family with arbitrarily many active coordinates.
The source's odd-degree geometric and elliptic continued-density inputs
remain outside the completed audit; the limit and its product corollary
do not require those inputs. No novelty, external-review or formal-verification
claim is made. General FC and unrestricted SFC remain open.

\section{A complete block limit, including inactive variables and markers}
Partition the active independent $X_i\sim\operatorname{Gamma}(\alpha_i,1)$,
$\alpha_i>0$, into nonempty blocks $B_1,\ldots,B_r$, $r\geq1$. Write
\[
 S_j=\sum_{i\in B_j}X_i,\quad A_j=\sum_{i\in B_j}\alpha_i,
 \quad F_0=\prod_j S_j^{d_j},\quad D=\sum_jd_j,
 \quad C_d=\prod_jd_j^{d_j},\qquad d_j\geq1.
\]
Let $T^{(j)}$ be the independent Dirichlet angular variables of these
blocks and abbreviate $dT=(d_1T^{(1)},\ldots,d_rT^{(r)})$.
Allow also any finite tuple of inactive independent variables
$Z_l\sim\operatorname{Gamma}(\beta_l,1)$, $\beta_l>0$, independent of
all active blocks.

The empty inactive tuple is allowed.
For a nonzero polynomial marker $H(X,Z)$ let $e=\deg_XH$ and let
$H_{[e]}$ be its part of active degree $e$, retaining its full dependence
on $Z$. The zero marker is identically zero.

\begin{theorem}\label{thm:block}
For arbitrary complex coefficients, $c\ne0$, homogeneous $G(X,Z)$ of
\emph{total} degree $D-1$, and $\deg R\leq D-2$, put
$F=cF_0+G+R$. Then
\[
 \lim_{m\to\infty}
 \frac{\E[H(X,Z)F(X,Z)^m]}{m^e c^m\prod_j(A_j)_{d_jm}}
 =\E_{T,Z}\left[H_{[e]}(dT,Z)
 \exp\!\left(\frac{G(dT,0)}{cC_d}\right)\right]. \tag{1}
\]
For $D=1$ the remainder is zero. Every shape may be any positive real.
A nonzero limit implies every sufficiently large marked moment is nonzero.
A zero limit makes no zero-sequence decision. With no inactive variables
and $H=1$, this is the probability-normalized pinned block-radial theorem.
\end{theorem}
\begin{proof}
Tilt only the active block radii by $F_0^m$, dividing by
$\E F_0^m=\prod_j(A_j)_{d_jm}$. Under this probability law the radii
are independent $\operatorname{Gamma}(d_jm+A_j,1)$, and their angles
retain their original Dirichlet laws. The inactive variables retain
exactly their fixed Gamma laws. Thus $S_j/m\to d_j$ in probability
(the variance is $d_j/m+A_j/m^2$), whereas $Z/m\to0$.

Choose a compact neighborhood $V$ of $d=(d_1,\ldots,d_r)$ inside
the positive orthant and fix $\delta>0$. Uniformly for $u\in V$,
$\sum_l v_l\leq\delta$, and all angles,
\[
 \frac{F(muT,mv)}{c\,m^D\prod u_j^{d_j}}
 =1+\frac{G(uT,v)}{cm\prod u_j^{d_j}}+O(m^{-2}).
\]
Its $m$th power is bounded by a constant independent of $m,u,v,T$.
For every fixed $Z$, put $v=Z/m$: uniformly for $u\in V$ and all
angles the power tends to
\[
 \exp\!\left(\frac{G(uT,0)}{c\prod u_j^{d_j}}\right),
 \qquad m^{-e}H(muT,Z)\longrightarrow H_{[e]}(uT,Z).
\]
Moreover, on $V$,
\[
 |m^{-e}H(muT,Z)|\leq C(1+\textstyle\sum_l Z_l)^L
\]
for some fixed integer $L$. This is integrable under the unchanged
inactive Gamma law. Conditional on fixed $T,Z$, concentration in $u$
gives the asserted integrand at $d$. Dominated convergence in $T,Z$
then proves the contribution from $u\in V$, $\sum Z_l/m\leq\delta$.

We now bound its complement in absolute value, including near zeros
of $F_0$, without dividing by $F_0$ there. Scale all radial variables
as $(S,Z)=m(u,v)$, and set
\[
 W=\sum_j u_j+\sum_l v_l,\qquad
 K_m(u,v,T)=e^{-W}|m^{-D}F(muT,mv)/c|.
\]
Uniformly on every compact set, including coordinate boundaries,
$K_m\to K_0=e^{-W}\prod u_j^{d_j}$.
The function $K_0$ has the unique global maximum
$M=e^{-D}C_d>0$ at $(u,v)=(d,0)$. It need not vanish when an
\emph{inactive} coordinate is zero. On a compact complement of
$V\times\{\sum v_l\leq\delta\}$ there is therefore a fixed $\rho<M$
such that $K_m\leq\rho$ for all sufficiently large $m$.
For the unbounded tail, the uniform polynomial bound gives, after
increasing a fixed threshold for $W$,
\[
 K_m\leq C e^{-W}(1+W)^D\leq\rho e^{-W/2}.
\]
The remaining radial weight is
$w(u,v)=\prod u_j^{A_j-1}\prod v_l^{\beta_l-1}$.
Its compact mass is finite even for shapes below one. If $h=\deg H$,
then $|m^{-e}H(muT,mv)|\leq C m^{h-e}(1+W)^h$.
Consequently the complementary scaled integral is bounded by a fixed
power of $m$ times $\rho^m$: on the tail use
$\int(1+W)^h e^{-mW/2}w(u,v)\,du\,dv=O(m^{-A-B})$,
where $A=\sum A_j$ and $B=\sum\beta_l$.
All angular measures are probability measures, including their faces.
Finally the scaled normalization is, by Stirling,
\[
 N_m:=\frac{\prod_j\Gamma(d_jm+A_j)}{m^{Dm+A+B}}
   =C_* M^m m^{-r/2-B}(1+o(1)),\qquad C_*>0.
\]
The additional fixed factors $\prod_l\Gamma(\beta_l)$ do not affect
the estimate. Dividing the complementary bound by $N_m$ makes it
exponentially small. This proves (1) for arbitrary complex lower
coefficients and every positive shape.
\end{proof}

\begin{corollary}[A single top monomial]\label{cor:monomialtop}
Suppose the highest ordinary homogeneous part of $F(Y)$ is a single
nonconstant monomial $cY^a$, with $a_i\geq0$ and $c\ne0$. Put
$I=\{i:a_i>0\}$, $D=\sum_i a_i$, $C_a=\prod_{i\in I}a_i^{a_i}$,
and let $G$ be the degree-$D-1$ homogeneous part of $F$.
For independent Gamma variables of arbitrary positive real shapes,
\[
 \frac{\E[F(Y)^m]}{c^m\prod_{i\in I}(\alpha_i)_{a_i m}}
 \longrightarrow \exp\!\left(\frac{G(a)}{cC_a}\right)\ne0.
\]
Thus every sufficiently large pure moment is nonzero, for arbitrary
complex strictly lower-degree terms in \emph{all} variables. For a
marker of active degree $e$, its normalized limit is this exponential
times $\E_Z H_{[e]}((a_i)_{i\in I},Z)$. If this amplitude is nonzero,
every sufficiently large marked moment is nonzero as well.
\end{corollary}
\begin{proof}
Take the positive-exponent variables as active singleton blocks and
the zero-exponent variables as inactive. The singleton angular laws
are point masses, so (1) gives both formulas.
\end{proof}
In particular, for the factorial functional
$\mathcal L(Y^b)=\prod_i b_i!$, all shapes equal one and
\[
 F=X^3+Q_{\leq2}(X,Y,Z)
 \quad\Longrightarrow\quad
 \frac{\mathcal L(F^m)}{(3m)!}
 \longrightarrow
 \exp\!\left(\frac{Q_2(3,0,0)}{27}\right)\ne0. \tag{1a}
\]
This allows arbitrary complex $XY,XZ,YZ,Y^2,Z^2$ couplings and
all affine terms. Likewise $X^2Y+Q_{\leq2}(X,Y,Z,W)$ has denominator
$(2m)!m!$ and limit $\exp(Q_2(2,1,0,0)/4)$.
The all-active example $XYZ+Q_{\leq2}$ remains included.
Inactive variables have ordinary fixed Gamma scale; their next-layer
couplings disappear from the exponential, while their marker moments remain.
For example, a unit-exponential inactive marker $Z-1$ has zero leading
amplitude, but for $F=X^3+XZ$ one has $\E[(Z-1)F]=1$.
Orthogonal Gaussian block lifts preserve the statement and original
markers after their actual angular averaging.

More generally, any proved nonzero unmarked block limit survives
adjoining inactive variables and arbitrary lower-degree couplings,
provided $G(X,0)$ remains in that proved angular family. In particular,
the source's block-separable next layer
\[
 G(X,0)=\sum_j d_j\xi_j\left(\prod_{l\ne j}S_l^{d_l}\right)
       S_j^{d_j-1-a_j}\widetilde Q_j(X^{(j)}),\qquad a_j\leq d_j-1,
\]
has $G(dT,0)/C_d=\sum_j\xi_jQ_j(T^{(j)})$. Its unit limit is the
product of the actual individual periods, independent of the inactive
couplings. Nonvanishing follows for the previously proved inputs, such as
\href{SIMPLEX_PENCILS_AND_AFFINE_LIMITS.tex}{integer-shape affine} and
\href{TWO_ANGULAR_COORDINATES.tex}{specified cubic} periods. The extension
does not audit the remaining geometric and elliptic inputs of EIT VI.

\section{An arbitrary number of coupled angular coordinates}
Let $U_1,\ldots,U_r$ be independent uniforms on $[0,1]$, $r\geq1$.
For $j\geq1$ define the entire functions
\[
 J_j(z)=\sum_{n\geq1}\frac{z^n}{n!n^j},\qquad J_0(z)=e^z-1.
\]
Their differential system is
\[
 (e^z)'=e^z,\quad 1'=0,\quad J_1'=(e^z-1)/z,
 \quad J_j'=J_{j-1}/z\quad(j\geq2).                  \tag{2}
\]
It has no finite pole except zero.

\begin{lemma}\label{lem:indep}
For every $k\geq0$, the functions $1,e^z,J_1,\ldots,J_k$ are
$E$-functions and linearly independent over $\C(z)$. At every
$\xi\in\Qb\setminus\{0\}$ their values are linearly independent over $\Qb$.
\end{lemma}
\begin{proof}
The normalized Taylor coefficients of $J_j$ are $1/n^j$ for $n\geq1$
and zero at zero. They are bounded rational numbers, with common
denominator dividing $\operatorname{lcm}(1,\ldots,N)^j\leq16^{jN}$.
For completeness, the quotient of the lcm through $2N$ by that through
$N$ divides $\binom{2N}{N}$: every newly appearing highest prime power
lies in $(N,2N]$ and contributes at least one to the binomial valuation.
Iterate this bound at powers of two and round up to get the displayed
exponential estimate. System (2) supplies the required differential
equations. Thus the arithmetic and analytic $E$-function conditions hold.

The functions $1,e^z$ are independent. If $J_1=a(z)e^z+b(z)$ with
rational $a,b$, differentiation and independence force $b'=-1/z$,
impossible for a rational function. Inductively, suppose the list through
$J_{k-1}$ is independent and express $J_k$ as a rational combination
of it. For $k\geq2$, comparison of the coefficient of $J_{k-1}$ after
differentiation gives $a_{k-1}'=1/z$, again impossible. A derivative of
a rational function has zero residue at every finite point.

Apply \href{https://arxiv.org/html/math/0405549}{Beukers, Theorem 3.2
and Lemma 3.1}, in the exact linear-lifting form checked in the
\href{LONG_EIT_AND_FC2_AUDIT.tex}{Long/EIT audit}. Every nonzero algebraic
$\xi$ is ordinary for (2). An algebraic linear relation of the values
would lift to a nonzero polynomial-coefficient functional relation with
the prescribed specialization, contradicting functional independence.
\end{proof}

\begin{theorem}\label{thm:product}
For every $r\geq1$ and $\xi\in\Qb\setminus\{0\}$,
\[
 \Phi_r(\xi):=\int_{[0,1]^r}e^{\xi\prod U_j}\,dU
       =\frac{J_{r-1}(\xi)}{\xi}
\]
is transcendental, hence nonzero. At zero it equals one.
\end{theorem}
\begin{proof}
Termwise integration gives
$\Phi_r(z)=\sum_{k\geq0}z^k/[k!(k+1)^r]=J_{r-1}(z)/z$,
with the removable value at zero. Lemma~\ref{lem:indep} proves the claim,
including $r=1$ where $J_0=e^z-1$.
\end{proof}
For $r\geq2$ the angular phase $\prod U_j$ is coupled; it is not a sum
of phases from the separate blocks. We do not multiply individual
periods to prove its nonvanishing.

\section{Exact decisions for every algebraic polynomial marker}
For $P(U)=\sum_a p_aU^a\in\Qb[U]$, put
\[
 W_P(t)=\sum_a p_a\prod_{i=1}^r\frac1{t+a_i+1},\qquad
 I_P(z)=\int_{[0,1]^r}P(U)e^{z\prod U_i}\,dU.
\]
Then $I_P(z)=\sum_{k\geq0}W_P(k)z^k/k!$.
All poles of the proper rational function $W_P$ are negative integers.
In its partial fractions write
$W_P(t)=\sum_{j,h}A_{j,h}/(t+j)^h$, $j\geq1$, $h\leq r$.
For signed Stirling numbers of the first kind $s(j,l)$,
\[
 \sum_{k\geq0}\frac{z^k}{k!(k+j)^h}
   =z^{-j}\sum_{l=1}^j s(j,l)J_{h-l}(z).             \tag{3}
\]
For negative indices use $J_{-b}=\theta^b(e^z-1)
=e^z\sum_l\left\{{b\atop l}\right\}z^l$, $b\geq1$.
Formula (3) follows by setting $n=k+j$ and expanding
$\fall nj=\sum_l s(j,l)n^l$. Extending the resulting sum down to $n=1$
adds zero, because $\fall nj=0$ for $1\leq n<j$.
Therefore, exactly,
\[
 I_P(z)=a_0(z)+a_e(z)e^z+\sum_{l=1}^{r-1}a_l(z)J_l(z),             \tag{4}
\]
with algebraic-coefficient Laurent polynomials $a_l(z)$.
At every nonzero algebraic $\xi$, Lemma~\ref{lem:indep} gives
\[
 I_P(\xi)=0\quad\Longleftrightarrow\quad
       a_0(\xi)=a_e(\xi)=a_1(\xi)=\cdots=a_{r-1}(\xi)=0.          \tag{5}
\]
If any coefficient other than $a_0(\xi)$ is nonzero, the period is
transcendental. Otherwise its value is the algebraic number $a_0(\xi)$.
At $\xi=0$ use the direct mass $W_P(0)$ instead of Laurent substitution.

This is an exact single-period zero test, not a test for a whole marked
power sequence. For example, at $r=2$, $\xi=1$,
\[
 P(u,v)=1-2u+u(1-u)v
\]
is a nonzero marker with $I_P(1)=0$, since
$P e^{uv}=\partial_u[u(1-u)e^{uv}]$. But $\E[P]=1/12\ne0$.
In a block limit, such a zero only removes the leading asymptotic
coefficient; lower corrections may survive.

\section{Higher-dimensional factorial and Gaussian families}
Take disjoint pairs of factorial variables $(X_j,Y_j)$ and write
$S_j=X_j+Y_j$. For $d_1\geq2$, $d_j\geq1$ otherwise, let
\[
 \begin{split}
 F_0&=\prod_{j=1}^r S_j^{d_j},\\
 G_\xi&=d_1\xi\,S_1^{d_1-2}
           \prod_{j=2}^rS_j^{d_j-1}\prod_{j=1}^r X_j,\\
 F&=F_0+G_\xi+K(X,Y,Z)+R(X,Y,Z),\qquad \deg R\leq D-2.
 \end{split}
 \tag{6}
\]
Here $Z$ is any finite independent tuple of positive-shape Gamma
variables, and $K$ is any homogeneous polynomial of total degree $D-1$
with $K(X,Y,0)=0$, with arbitrary complex coefficients. Thus every
strictly lower-degree coupling to inactive variables is allowed.
The variables named $X,Y$ in (6) are exponential/factorial variables,
not standard Gaussians. In the block limit each $U_j=X_j/S_j$ is uniform,
and $G_\xi(dT)/C_d=\xi\prod U_j$.

\begin{corollary}\label{cor:FC}
For every algebraic $\xi$, every polynomial (6), with arbitrary complex
coefficients in $K,R$, has all sufficiently large Gamma moments nonzero
(and factorial moments when all shapes are one).
For every complex $\xi$ and fixed algebraic inactive shapes, it
satisfies FC and has infinitely many nonzero moments. No all-start short SFC window is asserted.
\end{corollary}
\begin{proof}
For algebraic scales the unit limit is $\Phi_r(\xi)\ne0$.
For arbitrary complex scales an all-moment counterexample in this
finite-type family, defined over the algebraic number field of the
fixed inactive shapes, would specialize to an algebraic one. Explicitly, the ideal of all moment polynomials is finitely
generated; a proper such ideal has a $\Qb$-point by Nullstellensatz.
The leading part $F_0$ is fixed and nonzero. This contradicts the
algebraic-scale conclusion.

Powers preserve the displayed class: replace $d_j$ by $k d_j$ and
$G_\xi$ by $kF_0^{k-1}G_\xi$, which is exactly the new formula at the
same $\xi$. The new inactive next layer is $kF_0^{k-1}K$,
which still vanishes at $Z=0$; all other terms have degree at most $kD-2$.
An eventual zero tail would therefore contradict FC for a positive power.
\end{proof}

\begin{corollary}[Ordinary tails over every rational algebra]\label{cor:rings}
For each fixed $r,d$ and every starting index $N\geq1$, the universal
ideal generated by all moments of (6) at indices $m\geq N$ is the
\emph{ordinary unit ideal} in the rational polynomial ring of $\xi$ and
all coefficients of $K,R$ (for any fixed finite tuple of inactive
variables with positive rational shapes). Consequently, over every nonzero commutative
$\Q$-algebra such a normalized family has no eventual pure zero tail.
Its actual scalar normalization ideal for all marked periods is zero.
\end{corollary}
\begin{proof}
A proper tail ideal would have a complex point by Nullstellensatz,
contradicting Corollary~\ref{cor:FC}. A unit radical implies the ideal
itself contains one, so some finite ordinary Bezout combination of
the tail moments equals one. This identity persists under every base
change. In particular a scalar annihilating all pure moments eventually
is zero; the unit marker is among all actual markers. A unit overall
scalar multiplying $F$ gives the same conclusion. No explicit uniform
cutoff, short consecutive window or SFC assertion follows from this
existence argument.
\end{proof}
The same ordinary tail-unit and zero scalar-normalization conclusions
hold for the normalized single-top-monomial family with fixed positive
rational shapes, by Corollary~\ref{cor:monomialtop} and the identical
ordinary ideal argument.
For $r=2,d=(2,1)$ the concrete new coupled example is
\[
 (X_1+Y_1)^2(X_2+Y_2)+2\xi X_1X_2+R_{\leq1}.       \tag{7}
\]
It has four active factorial variables. With inactive variables, add
any quadratic $K$ vanishing at $Z=0$ and any affine $R$. For example,
$K=Z(X_1+Y_2)+Z^2$ gives genuine active--inactive coupling without
changing the limiting period.
The nonzero algebraic-scale limit is $J_1(\xi)/\xi$, not a product of
one-block periods.
For the subfamily $K=R=0$, the first two moments are
$\mu_1=12+2\xi$, $\mu_2=720+144\xi+16\xi^2$ and
$\mu_2-(8\xi+24)\mu_1=432$, an explicit ordinary witness in this
restricted subfamily. It is not a claimed cutoff for arbitrary $K,R$.

Each $X_j,Y_j$ can also be a sum of independent Gamma variables whose
shapes total one in each arm; the unit statement and arbitrary lower
terms still follow for every algebraic scale. The complex-scale
specialization statement applies when these fixed individual shapes are
algebraic. Standard Gaussian lifts use
$X_j=(x_{j1}^2+x_{j2}^2)/2$ and
$Y_j=(y_{j1}^2+y_{j2}^2)/2$, giving $4r$ active Gaussian coordinates.
Each inactive shape $k/2$, $k$ a positive integer, adds $k$ Gaussian
coordinates through its half squared norm.
Every original Gaussian marker is first averaged over these orthogonal
groups. Theorem~\ref{thm:block} retains the averaged marker's actual top
active-degree part, first taking the unchanged inactive Gamma
expectation; (4)--(5) then decide its leading cube period when the
resulting coefficients are algebraic.
An exact zero average is a whole zero-sequence identity; a zero leading
period alone is not.

\section{Executable scope and remaining gap}
\href{coupled_block_periods.py}{The exact implementation} verifies literal
block tops on a nonempty subset of the actual variables, active-degree
marker parts and exact inactive Gamma expectations for positive rational shapes,
reduces every algebraic cube marker by (3), and decides (5) at exact
algebraic scales. It exports the actual family (6), leading marker
periods, and the distinction between a zero period and a zero sequence.
\href{verify_coupled_block_periods.py}{Controls} compare independent
Gamma, Dirichlet and original Gaussian moments, the exponential Taylor
coefficients, Stokes zero controls, block-power identities and source
normalizations. Numerical limit samples are labeled diagnostics only.
The all-$r$ theorem uses the written independence and arithmetic lifting
proof, not a bounded rank or moment scan.

This adds coupled products on products of intervals, not every polynomial
of two angular coordinates on a triangle. Arbitrary coupled simplex
phases, the remaining EIT V geometric/elliptic audit, general FC and
unrestricted SFC remain open. There is no implication that all complex
scales have nonzero leading limits: already $\Phi_1(2\pi i)=0$.
\end{document}
